% from ega4-2-fr.tex
\oldpage[IV-2]{5}
\phantomsection
\label{chapter:IV2-fr}

\begin{center}
{\large CHAPITRE IV (suite)\par}
\vspace{1.5em}
{\Large\bfseries ÉTUDE LOCALE DES SCHÉMAS\\
ET DES MORPHISMES DE SCHÉMAS\par}
\end{center}

\setcounter{section}{1}
\section{Changement de base et platitude}
\label{section:IV.2-fr}

Ce paragraphe (contrairement au \S~6) ne fait appel qu'exceptionnellement à
des techniques noethériennes. Les n\textsuperscript{os}~1 et 2 ne sont guère
que des traductions de propriétés élémentaires de platitude en Algèbre
commutative (cf. Bourbaki, \emph{Alg. comm.}, chap.~I\textsuperscript{er}), et
sont incluses ici pour la commodité des références. Les numéros suivants sont
surtout consacrés à des propriétés de «~descente~» par les morphismes plats ou
fidèlement plats~: si $g:Y'\to Y$ est un tel morphisme, il s'agit de pouvoir
affirmer qu'une partie de $Y$, ou un $\mathcal O_Y$-Module, ou un morphisme
$X\to Y$, a une certaine propriété, lorsque l'on sait que son \emph{image
réciproque} par $g$ possède cette propriété. Nous nous bornons ici aux
propriétés ne faisant pas appel à la technique générale de la «~descente~»,
qui sera développée au chapitre V.

\setcounter{subsection}{0}
\subsection{Modules plats sur les préschémas}
\label{subsection:IV.2.1-fr}

\begin{env}[2.1.1]
\EGAReaderTarget{ega4.c04.statement.2-1-1}{IV.2.1.1-fr}
Soient $f:X\to Y$ un morphisme de préschémas, $\mathscr F$ un
$\mathcal O_X$-Module~; rappelons $(0_{\mathrm I},6.7.1)$ que $\mathscr F$ est
dit $f$-plat (ou $Y$-plat) \emph{en un point} $x\in X$ si $\mathscr F_x$ est
un $\mathcal O_{f(x)}$-module plat, $f$-plat (ou $Y$-plat) s'il est $f$-plat
en tout point $x\in X$, et enfin que le morphisme $f$ est dit plat \emph{au
point} $x\in X$ (resp. plat) si $\mathcal O_X$ est $f$-plat au point $x$
(resp. $f$-plat). Lorsque $f=1_X$, on dira simplement qu'un
$\mathcal O_X$-Module $\mathscr F$ est plat \emph{au point} $x$ (resp. plat)
s'il est $X$-plat en ce point (resp. en tout point $x\in X$), c'est-à-dire si
$\mathscr F_x$ est un $\mathcal O_x$-module plat (resp. s'il en est ainsi pour
tout $x\in X$). Rappelons que nous avons démontré (III, 1.4.15.1) la propriété
suivante~:
\end{env}

\begin{proposition}[2.1.2]
\label{IV.2.1.2-fr}
Soient $A$, $B$ deux anneaux, $\varphi:A\to B$ un homomorphisme d'anneaux,
$X=\operatorname{Spec}(B)$, $Y=\operatorname{Spec}(A)$, $f:X\to Y$ le
morphisme correspondant à $\varphi$, $M$ un $B$-module. Pour que
$\mathscr F=\widetilde M$ soit $f$-plat, il faut et il suffit que $M$ soit un
$A$-module plat.
\end{proposition}

